\documentclass{article}
\usepackage[T1]{fontenc}
\usepackage{lmodern}
\usepackage{graphicx}
\usepackage{amsmath,amssymb}
\usepackage[a4paper,margin=2cm]{geometry}
\usepackage[absolute,overlay]{textpos}
\setlength{\TPHorizModule}{1mm}
\setlength{\TPVertModule}{1mm}
\usepackage[indonesian]{babel}
\usepackage{hyperref}
\setlength{\emergencystretch}{2em}
% unit: Halaman 11

\begin{document}

% === Halaman 11 (python/native) ===

• Tujuan utama steganalisis adalah untuk membedakan apakah sebuah media 

mengandung pesan rahasia atau tidak.

• Steganalisis dianggap berhasil jika ia dapat menentukan apakah sebuah 

media mengandung pesan tersembunyi dengan peluang lebih tinggi 

daripada menerka secara acak.

• Selain tujuan utama di atas, terdapat beberapa tujuan minor steganalisis:

   - menentukan panjang pesan

   - menentukan tipe algoritma penyisipan 

   - kunci yang digunakan 


11

\end{document}
